\pdfoutput=1
\documentclass[11pt]{article}
\usepackage[T1]{fontenc}
\usepackage{lmodern,amsmath,amssymb,amsthm}
\usepackage[margin=1in]{geometry}
\usepackage[colorlinks=true,linkcolor=blue,citecolor=blue,urlcolor=blue]{hyperref}
\hypersetup{pdftitle={Homomorphisms to long odd cycles in time exponential in the order-to-length ratio},pdfauthor={Jason Davies}}
\newtheorem{theorem}{Theorem}
\newtheorem{lemma}[theorem]{Lemma}
\newtheorem{corollary}[theorem]{Corollary}
\newcommand{\Z}{\mathbb Z}
\title{Homomorphisms to long odd cycles in time exponential in the order-to-length ratio}
\author{Jason Davies}
\date{November 18, 2026}
\begin{document}
\maketitle

\begin{abstract}
Given an $N$-vertex graph $G$ and an odd integer $\ell\geq3$ in binary,
we decide $G\to C_\ell$ deterministically in $2^{O(N/\ell)}$ time up to
polynomial factors and in polynomial space. A homomorphism lifts along
a spanning tree to integer heights that change by one at every step.
Sampling these heights along a tree tour determines how many times each
edge difference crosses the modulus. Shortest paths then recover a
colouring. As a consequence, recognising an odd-cycle core takes
polynomial time on graphs of odd girth $\Omega(N/\log N)$, with no
restriction on even cycles.
\end{abstract}

\section{The target length as a parameter}
A homomorphism $G\to C_\ell$ labels each vertex by a residue modulo $\ell$
so that adjacent labels differ by $1$ or $-1$. How much does a long
target help? Fomin, Heggernes and Kratsch~\cite{fhk} enumerate the
preimage of a least-used colour and solve the remaining problem in
polynomial time. For odd $\ell\geq5$ their deterministic bound includes
the factor $\binom{N}{\lfloor N/\ell\rfloor}$. Carbonnel~\cite[Theorem~2]{carbonnel}
gives a randomised algorithm with exponential factor
\[
 \left[(\ell/2)^{1/\ell}
       \bigl(\ell/(\ell-1)\bigr)^{(\ell-1)/\ell}\right]^N .
\]
Both expressions give an exponent of order $(N/\ell)\log\ell$ when
$\ell$ grows and $N/\ell$ grows. We obtain an exponent linear in
$N/\ell$. For example, at $\ell=\Theta(N/\log N)$ the binomial bound is
quasipolynomial and our algorithm takes polynomial time.

\begin{theorem}\label{thm:long}
For any finite simple graph $G$ on $N$ vertices and odd $\ell\geq3$
given in binary, one can decide $G\to C_\ell$ and return a homomorphism
on a yes instance in $2^{O(N/\ell)}|I|^{O(1)}$ deterministic time and
polynomial space, where $|I|$ is the bit length of the input.
\end{theorem}

\section{Lifting a colouring along a tree}
Consider a connected graph with a rooted spanning tree $T$. Write $d_v$
for the depth of $v$. Rotate any homomorphism so the root has colour
zero. Lift it to integer heights $H_v$ by setting the root height to
zero and following the colour increments $\pm1$ along tree edges. Then
$H_v$ represents the colour of $v$ and $H_v\equiv d_v\bmod 2$.
Orient each graph edge arbitrarily. For every $uv$ there is an integer
$k_{uv}$ such that
\begin{equation}\label{eq:winding}
       H_v-H_u=k_{uv}\ell\pm1,\qquad
       k_{uv}\equiv d_v-d_u-1\pmod 2 .
\end{equation}
The parity follows because $\ell$ is odd. On tree edges $k_{uv}=0$.
The integers on the other edges record the multiple of $\ell$ traversed
by the lifted tree path between their endpoints, up to the final unit step.
For example, lift the identity map of $C_\ell$ along the path
$0,1,\ldots,\ell-1$. On the closing edge oriented from $0$ to $\ell-1$,
the height difference is $\ell-1$ and $k_{0,\ell-1}=1$.

\begin{lemma}\label{lem:solve}
Given one integer $k_{uv}$ of the prescribed parity for every oriented
edge, the existence of heights satisfying~\eqref{eq:winding} can be
tested in polynomial time in the input bit length. Feasible heights give
a homomorphism.
\end{lemma}
\begin{proof}
Use the inequalities $H_v-H_u\leq k_{uv}\ell+1$ and
$H_u-H_v\leq 1-k_{uv}\ell$. Represent each as a directed arc from the
subtracted variable to the other, with its right-hand side as weight.
A negative cycle makes the inequalities inconsistent. Otherwise let
$H_v$ be the shortest distance from the root in this digraph. Every
vertex is reachable because each graph edge supplies both arcs.
The distances satisfy both inequalities. Each arc weight has the parity
of the depth difference of its endpoints, so every root-to-$v$ walk has
weight congruent to $d_v$ modulo two. Thus $H_v-H_u$ cannot equal the
midpoint $k_{uv}\ell$, and the only remaining values are the two in
\eqref{eq:winding}. Reducing the distances modulo $\ell$ gives the
colouring. Bellman--Ford computes the distances or detects a negative cycle.
\end{proof}

\section{Sampling a tree tour}
\begin{proof}[Proof of Theorem~\ref{thm:long}]
First consider a connected component of order $m$. Bipartite components
map to an edge. If $\ell>m$, a nonbipartite component cannot map to
$C_\ell$, since its odd cycle would give an odd closed walk shorter than
$\ell$ in $C_\ell$. For $\ell=3$ one may enumerate the two signs on each
tree edge and check the remaining edges in time $2^{m-1}|I|^{O(1)}$.
Hence assume $5\leq\ell\leq m$.

Put $s=\lfloor(\ell-1)/3\rfloor$. Start at the root and tour $T$ by
traversing each edge once in each direction. Write $v_t$ for the vertex
at time $t$, where $0\leq t\leq2(m-1)$. Sample times $0,s,2s,\ldots$
and, for a putative lift, write
\[
 b_i=\left\lfloor H_{v_{is}}/s\right\rfloor .
\]
We know $b_0=0$. Consecutive sampled heights are at distance at most $s$,
so $b_{i+1}-b_i\in\{-1,0,1\}$, also for negative heights.
Enumerate all such sequences. For each vertex $v$ fix its first occurrence
at time $t$ and let $i=\lfloor t/s\rfloor$. At most $s$ steps separate
it from the preceding sample. Thus the true lift belongs to the interval
\[
             J_v=[s b_i-s,\ s b_i+2s-1].
\]
Fixing a single occurrence per vertex suffices even when the tour
revisits it. The last incomplete block uses its preceding sample as well.

For an oriented edge $uv$ with $J_u=[a,b]$ and $J_v=[c,d]$, equation
\eqref{eq:winding} requires $k_{uv}\ell\in[c-b-1,d-a+1]$. This interval
has width $6s<2\ell$. Multiples $k\ell$ of the prescribed parity of $k$
are spaced $2\ell$ apart. There is consequently at most one candidate,
which integer division locates. Reject the sequence if an edge has no
candidate. Otherwise apply Lemma~\ref{lem:solve} to the candidates for
all edges. The intervals have served their purpose once the integers
are determined.

Any homomorphism has a lift whose sampled sequence is enumerated.
On that branch the unique candidates are its own $k_{uv}$, so the
inequalities are feasible. In the other direction every successful branch
returns a homomorphism by the lemma. There are at most
$3^{\lfloor2(m-1)/s\rfloor}$ branches. Each sequence and its constraint
digraph have polynomial size, and the heights and arc weights require
only polynomially many bits. Processing branches in turn proves the time
and space bounds for a component. Solve components independently and
combine their colourings. Summing their times gives the asserted overall
bound, also for isolated vertices and the empty graph.
\end{proof}

\section{Recognising odd-cycle cores}
\begin{corollary}
Let $G$ have $N$ vertices and finite odd girth $h$. Odd-cycle core
recognition takes $2^{O(N/h)}|I|^{O(1)}$ time. In particular, it takes
polynomial time for
$h=\Omega(N/\log N)$.
\end{corollary}
\begin{proof}
Compute $h$ as the minimum distance from $(v,0)$ to $(v,1)$ over all
vertices $v$ in the bipartite double cover. This minimises the length
of an odd closed walk, which contains an odd cycle. The graph maps to
$C_h$ exactly when its core is an odd cycle. Indeed $C_h$ embeds in the
graph, every odd cycle is a core, and homomorphisms cannot increase odd
girth. Apply Theorem~\ref{thm:long} with $\ell=h$. Bipartite inputs can
be rejected separately.
\end{proof}

The corollary permits arbitrarily short even cycles. For homomorphism to
an independently supplied target it is the length of that target which
enters Theorem~\ref{thm:long}. A companion paper~\cite{structural} gives
hardness for source pathwidth, treewidth and feedback vertex number at
ordinary girth $\Omega(N^{1-\varepsilon})$ for every fixed
$0<\varepsilon<1$.

\paragraph{Acknowledgement.}
Mathematical development, checking and drafting used AI assistance under
human direction.
\begin{thebibliography}{9}
\bibitem{carbonnel}
Cl\'ement Carbonnel. Faster algorithms for graph homomorphism via tractable
constraint satisfaction. \href{https://arxiv.org/abs/2602.23000v2}{arXiv:2602.23000v2}, 2026.
\bibitem{structural}
Jason Davies. Parameterised complexity of graph homomorphism to odd cycles.
Companion manuscript, 2026. Distributed as \texttt{odd-cycle-homomorphism.pdf}.
\bibitem{fhk}
Fedor V. Fomin, Pinar Heggernes and Dieter Kratsch.
Exact algorithms for graph homomorphisms.
\emph{Theory of Computing Systems} 41:381--393, 2007.
\href{https://doi.org/10.1007/s00224-007-2007-x}{doi:10.1007/s00224-007-2007-x}.
\end{thebibliography}
\end{document}
